% =====================================================================
%  thmsmith Stage -1 proposal skeleton.
%  Written automatically by the thmsmith TS pipeline before Stage -1 runs.
%  Locked parameters: qid=pid\_bunching\_heaping\_phase, specialization=v1
%
%  HOW TO USE THIS FILE
%  --------------------
%  - The preamble (everything above the `% --- BEGIN CONTENT ---` marker)
%    and the `\section{...}` headers below are fixed by the pipeline.
%    Do NOT rewrite or rename them; the Step 3 reviewer subagent and the
%    downstream Stage 0 / Stage 1 prompts grep for these section titles.
%  - Each section contains exactly one `% TODO(stage-1 ...): ...`
%    placeholder. Replace each TODO with the content described for that
%    section in `CausalSmith/tools/src/discovery/prompts/stage_neg1_2_draft.txt`
%    ("REQUIRED OUTPUT FORMAT (Stage -1 proposal .tex)").
%  - The `\paragraph{Locked parameters.}` block in the metadata block
%    must pin the chosen `(qid, specialization, p, assignment, target,
%    regression)` precisely. If the proposal maps to the Q1 pool-mode,
%    reproduce that block verbatim from
%    `CausalSmith/doc/panel_formula_question_generator_note_with_common_prompts.tex`
%    (Q2..Q6 templates have been retired). Otherwise (free-form
%    `(qid, specialization)` — the default), define each parameter
%    explicitly here (no implicit conventions). Stage 0 quotes this block;
%    do not rephrase it after locking.
%  - Removing a section header, renaming a macro, or rewriting the
%    preamble is a regression: the file must still match this skeleton
%    after you finish.
% =====================================================================

\documentclass[11pt]{article}

\usepackage{amsmath,amssymb,amsthm,mathtools}
\usepackage{enumitem}
\usepackage[colorlinks=true,linkcolor=blue,citecolor=blue,urlcolor=blue]{hyperref}
\usepackage{natbib}

\newcommand{\E}{\mathbb{E}}
\renewcommand{\Pr}{\mathbb{P}}
\newcommand{\one}{\mathbf{1}}
\newcommand{\transpose}{\mathsf{T}}
\newcommand{\calH}{\mathcal{H}}
\newcommand{\calF}{\mathcal{F}}
\newcommand{\calR}{\mathcal{R}}
\newcommand{\R}{\mathbb{R}}
\newcommand{\plim}{\operatorname*{plim}}

\theoremstyle{definition}
\newtheorem{definition}{Definition}
\newtheorem{assumption}{Assumption}
\newtheorem{example}{Example}

\theoremstyle{plain}
\newtheorem{theorem}{Theorem}
\newtheorem{conjecture}{Conjecture}
\newtheorem{proposition}{Proposition}
\newtheorem{lemma}{Lemma}
\newtheorem{corollary}{Corollary}

\theoremstyle{remark}
\newtheorem{remark}{Remark}

\title{thmsmith Stage -1 proposal: \texttt{pid\_bunching\_heaping\_phase} / \texttt{v1}%
  \\ \large\textit{Sharp aliasing frontiers for bunched running variables observed only after deterministic heaping}}
\author{thmsmith Stage -1 proposer}
\date{\today}

\begin{document}
\maketitle

% --- BEGIN CONTENT ---  (everything above is fixed by the pipeline)

\paragraph{Locked parameters.}
\texttt{qid = pid\_bunching\_heaping\_phase; specialization = v1; observable distribution P = law of the rounded score H=c(X), an outcome-free bunching histogram on a finite latent refinement G=\{1,\ldots,K\}; parameter theta = excess latent bunching mass theta(p)=b^\transpose p relative to a fixed counterfactual stencil b supported on the excluded bunching window B and neighboring source cells S; identifying restrictions = p\in\Delta_K, Ap=q(P) for the known deterministic coarsening matrix A induced by c, fixed zero-mass normalization \one^\transpose p=1, and the shape class C_r(M) of side-window finite-difference inequalities of order r with curvature bound M; bounding functional = [\underline\theta(P),\overline\theta(P)] where \underline\theta(P)=\inf\{b^\transpose p:p\in\Delta_K,Ap=q(P),p\in C_r(M)\} and \overline\theta(P)=\sup\{b^\transpose p:p\in\Delta_K,Ap=q(P),p\in C_r(M)\}.}

\begin{abstract}
This question asks what a bunching design identifies when the latent running variable is observed only after deterministic rounding or heaping. The proposed focal object is the heap-aliasing frontier, the sharp interval for excess latent bunching mass over all latent histograms that induce the same rounded histogram. A setup lemma records the algebraic support interval and the feasibility condition for signed-nullspace certificates. The flagship conjecture gives a target-sensitive phase threshold in the heap-grid width and the excluded-window geometry: once a heap cell straddles the estimand boundary with a target-moving latent refinement, no estimator can recover the mass without additional shape restrictions. The result would separate genuine bunching information from artifacts of reported-score aggregation in notch, kink, and donut-style empirical designs.
\end{abstract}

\section{Introduction}
Reported running variables in tax, benefit, and administrative designs are often rounded at the same locations where bunching is measured. The main claim is that deterministic heaping creates an aliasing frontier: two latent score densities can agree on every rounded histogram cell and still imply different excess bunching mass.

The closest comparators are \citet[Sec. 2--3]{Song2024GeneralBunching}, which studies general bunching identification through counterfactual density restrictions, and \citet[Sec. 2]{Dong2015RoundingRD}, which shows that rounded running variables can invalidate standard regression discontinuity estimation. This proposal differs by making the coarsening operator itself the object of partial identification and by requiring an explicit signed-nullspace certificate rather than a deconvolution or polynomial extrapolation formula. The kernel is a PartialID theorem for deterministic coarsening in bunching designs.

The direct consumer is empirical bunching at tax notches such as \citet{KlevenWaseem2013Notches}, where reported income can be heaped and the bunching window is chosen by the researcher. A reader who accepts the result would report whether the window-grid alignment has any identifying content before estimating an elasticity from excess mass.

\begin{itemize}[leftmargin=*]
\item Theorem 1: The heap-aliasing frontier is the support interval of the latent feasible set, and signed nullspace directions certify nonidentification when the displayed baseline already matches the observed rounded histogram. Gap: \citet{ManskiTamer2002IntervalData} and \citet{BeresteanuMolchanovMolinari2011ConvexMoment} give the generic convex-support machinery; this lemma only fixes the bunching object used by the conjectures. Consumer: rounded-score bunching designs in \citet{KlevenWaseem2013Notches}.
\item Conjecture 1: In the equal-grid design, target-sensitive heap-null directions give the phase boundary for nontrivial aliasing of excess mass under bounded side curvature. Gap: \citet[p. 268]{BarrecaLindoWaddell2016HeapingRD} treats heaping as bias and \citet[pp. 951--979, Sec. 2]{BoschDekkerStrohmaier2020Window} treats the window as a tuning choice, but neither gives an informative/uninformative frontier. Consumer: bunching-window selection in tax and transfer designs.
\item Conjecture 2: Endpoint pruning is sharp only under an endpoint-cell dominance condition; bounded curvature alone leaves signed heap-null witnesses. Gap: \citet{BertanhaMcCallumPayneSeegert2022Stata} emphasizes unverifiable distributional assumptions, while \citet{ManskiTamer2002IntervalData} permits interval interiors without the bunching-local endpoint failure. Consumer: software and applied checks that replace rounded scores by interval endpoints.
\end{itemize}

\section{Related work}
\citet{Saez2010Bunching} introduced kink bunching estimators that translate excess mass into elasticities under a smooth counterfactual density. \citet{ChettyFriedmanOlsenPistaferri2011Adjustment} use bunching around tax kinks to study adjustment costs and motivate finite windows around policy notches and kinks. \citet{KlevenWaseem2013Notches} use excess bunching and missing mass around notches; the present question asks whether the excess mass is identified when reported scores are heaped. \citet{BertanhaMcCallumSeegert2021BetterBunching} show that flexible heterogeneity weakens kink identification and makes assumptions central. \citet{BertanhaMcCallumPayneSeegert2022Stata} implement nonparametric and semiparametric bunching methods and stress that assumptions on the unobserved distribution cannot be learned from the histogram alone. \citet{Song2024GeneralBunching} gives recent point and partial identification results for general bunching designs, with analyticity central to point identification.

\citet{Dong2015RoundingRD} studies RD with a rounded running variable and shows that standard estimation can be inconsistent. \citet{DaveziesLeBarbanchon2017RDMeasurementError} show that measurement error smooths discontinuities and can destroy RD identification without auxiliary structure. \citet{DongKolesar2023IgnoreMeasurementError} give conditions under which running-variable measurement error can be ignored in RD. \citet{BarrecaLindoWaddell2016HeapingRD} show that heaping can bias RD estimates and that standard diagnostics may miss it. \citet{McCrary2008DensityTest} supplies the density-continuity benchmark for manipulation tests, but deterministic heaping can preserve rounded counts while changing latent excess mass. \citet{KolesarRothe2018DiscreteRD} study inference with discrete running variables, which is adjacent but does not give a latent heap-aliasing frontier. \citet{NoackRothe2023DonutRD} analyze donut trimming; here trimming can remove the constraints needed to rule out aliasing. \citet{BoschDekkerStrohmaier2020Window} propose a data-driven bunching-window procedure; the present phase frontier asks when that window-grid geometry contains identifying information.

\citet{ManskiTamer2002IntervalData} give partial-identification bounds with interval covariates and no restrictions on the unobserved values inside intervals. \citet{BeresteanuMolchanovMolinari2011ConvexMoment} characterize sharp identification regions for interval and random-set models. These papers are the closest coarsening comparators, but the proposed result targets a local excess-mass functional and a signed nullspace witness tied to the deterministic heap map.

The CausalSmith catalogue records no exact prior bunching/heaping proposal in \texttt{doc/API.md}. The closest in-repo prior art is \texttt{pid\_rdd\_manipulation\_bounded}, whose headline uses observable cutoff-side weak limits under bounded manipulation, \texttt{pid\_rdd\_flow\_manipulation}, whose rejected many-window angle failed to compute a global endpoint envelope, \texttt{pid\_interval\_censored\_adoption}, whose latent-refinement witness concerns DiD event-time, and \texttt{pid\_multivariate\_interval\_regression}, whose rejected endpoint-pruning logic warns that interval endpoints need a dominance condition.

\section{Formal setup}
Let $X$ be a latent one-dimensional running variable supported on a finite refinement $G=\{1,\ldots,K\}$. The researcher observes only $H=c(X)\in\{1,\ldots,J\}$, where $c$ is known and deterministic. Let $A\in\{0,1\}^{J\times K}$ be the coarsening matrix with $A_{jk}=1$ if $c(k)=j$ and zero otherwise. The observed distribution is $P$, summarized by the rounded histogram $q(P)\in\Delta_J$, where $q_j(P)=\Pr(H=j)$.

The latent histogram is $p\in\Delta_K$. The bunching window is $B\subset G$, the side cells are $S\subset G\setminus B$, and the estimand is a fixed linear excess-mass functional
\[
  \theta(p)=b^\transpose p
\]
where $b_k=1$ on bunching cells and $b_k$ includes the chosen counterfactual-density stencil on side cells. The shape class $C_r(M)$ consists of latent histograms satisfying finite-difference inequalities $|D^r p|_\infty\le M$ on $S\cup B$ and any sign or monotonicity restrictions explicitly invoked below.

The identified set is
\[
  \Theta_I(P)=\{b^\transpose p:p\in\Delta_K,\ Ap=q(P),\ p\in C_r(M)\}.
\]
The heap-aliasing frontier is the ordered pair
\[
  \mathfrak A(P;A,b,C_r(M))=(\underline\theta(P),\overline\theta(P)).
\]
For a heap cell $F_j=c^{-1}(j)$ with ordered endpoints $\ell_j=\min F_j$ and $u_j=\max F_j$, define the endpoint replacement operator $E_j:\Delta_K\to\Delta_K$ by leaving all cells outside $F_j$ unchanged and replacing the within-cell vector by
\[
  (E_jp)_{\ell_j}=\lambda_j(p)\sum_{k\in F_j}p_k,\quad
  (E_jp)_{u_j}=(1-\lambda_j(p))\sum_{k\in F_j}p_k,\quad
  (E_jp)_k=0\quad(k\in F_j\setminus\{\ell_j,u_j\}),
\]
where $\lambda_j(p)\in[0,1]$ is the endpoint interpolation weight used by the endpoint-pruning rule. Let $E$ denote the composition of the $E_j$ over heap cells intersecting $B$, and let the endpoint-pruned frontier be
\[
  \mathfrak A_E(P)=
  \left[
  \inf_{\substack{p\in\Delta_K\cap C_r(M):Ap=q(P)}} b^\transpose Ep,\ 
  \sup_{\substack{p\in\Delta_K\cap C_r(M):Ap=q(P)}} b^\transpose Ep
  \right].
\]
\begin{quote}
\texttt{qid = pid\_bunching\_heaping\_phase; specialization = v1; observable distribution P = law of the rounded score H=c(X), an outcome-free bunching histogram on a finite latent refinement G=\{1,\ldots,K\}; parameter theta = excess latent bunching mass theta(p)=b'\!p relative to a fixed counterfactual stencil b supported on the excluded bunching window B and neighboring source cells S; identifying restrictions = p in Delta\_K, Ap=q(P) for the known deterministic coarsening matrix A induced by c, fixed zero-mass normalization 1'\!p=1, and the shape class C\_r(M) of side-window finite-difference inequalities of order r with curvature bound M; bounding functional = [theta\_lower(P),theta\_upper(P)] where theta\_lower(P)=inf\{b'\!p:p in Delta\_K, Ap=q(P), p in C\_r(M)\} and theta\_upper(P)=sup\{b'\!p:p in Delta\_K, Ap=q(P), p in C\_r(M)\}.}
\end{quote}

\section{Assumptions}
\begin{assumption}[Known deterministic heap map]\label{ass:heap}
The coarsening map $c:G\to\{1,\ldots,J\}$ is known, deterministic, and induces the matrix $A$ used in the definition of $\Theta_I(P)$.
\end{assumption}

\begin{assumption}[Finite latent refinement]\label{ass:finite}
The latent histogram $p$ lies in the simplex $\Delta_K$ on a fixed finite refinement $G=\{1,\ldots,K\}$, and the rounded histogram $q(P)$ is observed without sampling error at the identification stage.
\end{assumption}

\begin{assumption}[Fixed bunching stencil]\label{ass:stencil}
The bunching window $B$, side cells $S$, and linear stencil $b\in\R^K$ are fixed before observing $P$, with $b$ not proportional to any row of $A$ on the cells intersecting $B\cup S$.
\end{assumption}

\begin{assumption}[Shape class]\label{ass:shape}
For a fixed order $r\ge 1$ and constant $M\ge0$, admissible histograms satisfy $p\in C_r(M)$, where $C_r(M)$ is the intersection of $\Delta_K$ with the finite-difference inequalities $|D^r p|_\infty\le M$ on $S\cup B$ and any declared monotonicity or endpoint-dominance inequalities.
\end{assumption}

\begin{assumption}[Interior feasibility for certificate moves]\label{ass:interior}
When a signed direction $h$ is used as a certificate, there is an $\varepsilon>0$ such that $p_0\pm\varepsilon h\in\Delta_K\cap C_r(M)$ for the displayed baseline $p_0$.
\end{assumption}

\begin{assumption}[Side-stencil completeness]\label{ass:complete}
In the equal-grid phase design, every feasible target-moving first-order perturbation supported on $c^{-1}c(B)\cup S$ is represented in the local heap-null tangent space $L_b(A,B,r,M)$ defined in Conjecture \ref{conj:phase}. Equivalently, if $h\in\ker(A)$ is a feasible local tangent direction and $b^\transpose h\ne0$, then its projection onto $c^{-1}c(B)\cup S$ has a representative in $L_b(A,B,r,M)$ with nonzero $b$-contrast.
\end{assumption}

\begin{assumption}[Endpoint-cell dominance]\label{ass:dominance}
For every heap cell $F_j=c^{-1}(j)$ intersecting $B$ and every signed vector $u$ with $\one^\transpose u=0$, $\operatorname{supp}(u)\subseteq F_j$, and $p+\tau u\in C_r(M)$ for small $\tau>0$, the endpoint replacement rule $E_j$ satisfies
\[
  b^\transpose u\le 0\quad\Longrightarrow\quad b^\transpose(E_j(p+\tau u)-E_jp)\le 0,
\]
with the reversed inequality also holding after replacing $u$ by $-u$.
\end{assumption}

\section{Main results}
\begin{theorem}[Heap-aliasing frontier and signed-nullspace certificate]\label{thm:frontier}
Under Assumptions \ref{ass:heap}--\ref{ass:shape},
\[
  \Theta_I(P)=\left[\inf_{p\in\Delta_K\cap C_r(M):Ap=q(P)}b^\transpose p,\ 
  \sup_{p\in\Delta_K\cap C_r(M):Ap=q(P)}b^\transpose p\right].
\]
If there exist $p_0\in\Delta_K\cap C_r(M)$ with $Ap_0=q(P)$ and $h\in\ker(A)$ satisfying Assumption \ref{ass:interior} and $b^\transpose h\ne0$, then $\underline\theta(P)<\overline\theta(P)$.
\end{theorem}
The proof is the support-function identity for the image of a convex feasible set under a linear map. The certificate part follows from $A(p_0+\varepsilon h)=A(p_0-\varepsilon h)$ and $\theta(p_0+\varepsilon h)-\theta(p_0-\varepsilon h)=2\varepsilon b^\transpose h$.

\begin{conjecture}[Equal-grid heap/window phase threshold (NOT YET PROVED)]\label{conj:phase}
Assume Assumptions \ref{ass:heap}--\ref{ass:shape} and \ref{ass:complete} in an equal-latent-grid design where every heap cell contains $g$ consecutive latent cells and $B$ is a consecutive excluded window. Define the straddling count
\[
  \kappa(A,B)=\#\{j: c^{-1}(j)\cap B\ne\emptyset,\ c^{-1}(j)\cap (G\setminus B)\ne\emptyset\}
\]
and the local heap-null tangent space
\[
  L_b(A,B,r,M)=\{h\in\ker(A):\operatorname{supp}(h)\subseteq c^{-1}c(B)\cup S,\ D^r h=0\text{ on side cells where }M=0\}.
\]
The target-sensitive free-within-heap rank is the genuine linear rank
\[
  \rho_b(A,B,r,M)=\dim L_b(A,B,r,M)-\dim\bigl(L_b(A,B,r,M)\cap\ker(b^\transpose)\bigr).
\]
If $\kappa(A,B)=0$ or $\rho_b(A,B,r,M)=0$, then Assumption \ref{ass:complete} rules out every local heap-null perturbation that can move the target, so the side-stencil restrictions point identify the bunching mass within the equal-grid phase class. If $\kappa(A,B)\ge1$ and $\rho_b(A,B,r,M)\ge1$, then there is an open set of feasible rounded histograms $P$ for which $\underline\theta(P)<\overline\theta(P)$.
\end{conjecture}

\begin{conjecture}[Endpoint pruning requires endpoint dominance (NOT YET PROVED)]\label{conj:endpoint}
Under Assumptions \ref{ass:heap}--\ref{ass:shape}, endpoint pruning is the assertion $\mathfrak A_E(P)=\mathfrak A(P;A,b,C_r(M))$. This equality can hold uniformly over rounded histograms only if Assumption \ref{ass:dominance} holds. If Assumption \ref{ass:dominance} fails and a straddling heap cell contains at least three latent subcells, then bounded curvature alone admits a signed $h\in\ker(A)$ with $b^\transpose h\ne0$ and $\mathfrak A_E(P)\subsetneq \mathfrak A(P;A,b,C_r(M))$ for an open set of feasible $P$.
\end{conjecture}

\section{Sanity-check corollaries}
\begin{corollary}[No-heaping reduction]\label{cor:noheap}
Under Assumptions \ref{ass:heap}--\ref{ass:shape}, if $A$ is the identity matrix, then $\Theta_I(P)=\{\theta(q(P))\}$.
\end{corollary}
This collapses to ordinary observed-score bunching because $Ap=q(P)$ pins $p=q(P)$.

\begin{corollary}[Heap cells aligned with the window]\label{cor:aligned}
If every heap cell is either contained in $B$ or disjoint from $B$, and the stencil $b$ is constant within every heap cell outside the side-window extrapolation component, then heap aggregation alone does not create aliasing for the raw bunching count $\sum_{k\in B}p_k$.
\end{corollary}
This reduction follows because the bunching indicator lies in the row span of $A$.

\paragraph{Exhibit 9.1: signed heap-null certificate and phase object.}
Let $G=\{1,2,3,4\}$, $c(1)=c(2)=1$, $c(3)=c(4)=2$, $B=\{2,3\}$, and $b=(0,1,1,0)^\transpose$. Then
\[
  A=\begin{pmatrix}1&1&0&0\\0&0&1&1\end{pmatrix},\quad
  h=(1,-1,1,-1)^\transpose,\quad Ah=0,\quad b^\transpose h=0.
\]
The raw count is not changed by this symmetric direction, so the internal-consistency check rejects it as a certificate for the chosen $b$. For the excess-mass stencil $b^\star=(-1,1,1,-1)^\transpose$, the same primitive design gives $b^{\star\transpose}h=-4$. With $p_0=(1/4,1/4,1/4,1/4)$ and $0<\varepsilon<1/4$, $p_0\pm\varepsilon h$ are valid histograms with identical rounded histogram $(1/2,1/2)$ and different excess mass. Here $\kappa(A,B)=2$ and $\rho_{b^\star}(A,B,1,M)>0$ for any $M$ large enough to contain the two perturbations, so the exhibit computes the focal certificate from $A$, $B$, $b^\star$, and $p_0$ rather than assuming it.

\paragraph{Exhibit 9.2: three-subcell endpoint-pruning failure.}
Let one straddling heap cell be $F=\{1,2,3\}$ with rounded total $q_F=3/10$, and let two neighboring side cells have fixed masses $p_0(4)=p_0(5)=7/20$. Take
\[
  A_F=(1\ 1\ 1),\quad b_F=(-1,2,-1)^\transpose,\quad
  p_0=(1/10,1/10,1/10,7/20,7/20).
\]
The endpoint replacement rule keeps only cells $1$ and $3$ and assigns the heap mass symmetrically, giving $E_Fp_0=(3/20,0,3/20,7/20,7/20)$ and $b_F^\transpose E_Fp_0=-3/10$. The feasible interior perturbation $h_F=(1,-2,1,0,0)^\transpose$ has $A_Fh_F=0$ and $b_F^\transpose h_F=-6$. For any $0<\varepsilon<1/20$, $p_0\pm\varepsilon h_F$ have the same rounded histogram, remain nonnegative, and satisfy the same first-difference bound after increasing $M$ by at most $3\varepsilon$. Their target values differ by $12\varepsilon$, while the endpoint rule returns the same value for both. The endpoint frontier therefore misses feasible latent targets even though every symbol in the calculation is determined by $A_F$, $q_F$, $b_F$, and the displayed side masses; the failure is not a normalization artifact.

\section{Proof-sketch route}
Theorem \ref{thm:frontier} follows by writing the latent feasible set as $K(P)=\Delta_K\cap C_r(M)\cap\{p:Ap=q(P)\}$ and applying the support interval of a linear functional on a convex set. Conjecture \ref{conj:phase} should be proved by the chain
\[
\text{phase headline}\leftarrow[\text{unfold }\S6]\leftarrow[\text{substitute }\S7]\leftarrow
\text{row-span separation of }b\leftarrow
\text{local basis for }\ker(A)\text{ on straddling cells}\leftarrow
\text{finite-difference cone tangent calculation}\leftarrow
\text{open-feasibility perturbation lemma}.
\]
The substantive steps are the row-span separation, the local nullspace basis for $L_b(A,B,r,M)$, the tangent-cone calculation under Assumption \ref{ass:complete}, and the open-feasibility perturbation, so $k=4$ for the flagship conjecture. Conjecture \ref{conj:endpoint} should be proved by comparing the endpoint-pruned support interval $\mathfrak A_E(P)$ to the support function of $K(P)$, then extending Exhibit 9.2 from one three-subcell heap to all straddling cells through a direct-sum decomposition of heap-null tangent cones.

\section{Honest scope flags}
Theorem \ref{thm:frontier} is routine convex analysis and is included to fix the object. Conjectures \ref{conj:phase} and \ref{conj:endpoint} are not yet proved. Assumption \ref{ass:complete} is a substantive side-stencil condition; without it, Conjecture \ref{conj:phase} should be read only as a local no-aliasing statement on $L_b(A,B,r,M)$. The proposal deliberately does not claim deconvolution, elasticity identification, or asymptotic inference.

\section{Iteration trail}
\begin{itemize}[leftmargin=*]
\item v1 (cold-start) — initial draft around the heap-aliasing frontier, signed-nullspace certificate, and equal-grid phase conjecture; prior verdict: none.
\item v2 (revise) — corrected the window-paper citation, demoted the generic support result to a setup lemma, added $Ap_0=q(P)$, made the phase rank target-sensitive, formalized endpoint dominance, and added a three-subcell endpoint-pruning exhibit; prior verdict: REVISE.
\item v3 (revise) — redefined $\rho_b$ as a linear rank, added side-stencil completeness, defined the endpoint replacement operator and endpoint-pruned frontier, and restated Conjectures 1--2 against those objects; prior verdict: REVISE.
\end{itemize}

\section*{Bibliography}
\begin{thebibliography}{99}
\bibitem[Barreca, Lindo, and Waddell(2016)]{BarrecaLindoWaddell2016HeapingRD}
Barreca, A. I., J. M. Lindo, and G. R. Waddell (2016).
Heaping-induced bias in regression-discontinuity designs.
\emph{Economic Inquiry} 54(1), 268--293.

\bibitem[Beresteanu, Molchanov, and Molinari(2011)]{BeresteanuMolchanovMolinari2011ConvexMoment}
Beresteanu, A., I. Molchanov, and F. Molinari (2011).
Sharp identification regions in models with convex moment predictions.
\emph{Econometrica} 79(6), 1785--1821.

\bibitem[Bertanha, McCallum, and Seegert(2021)]{BertanhaMcCallumSeegert2021BetterBunching}
Bertanha, M., A. H. McCallum, and N. Seegert (2021).
Better bunching, nicer notching.
Finance and Economics Discussion Series 2021-002.

\bibitem[Bertanha, McCallum, Payne, and Seegert(2022)]{BertanhaMcCallumPayneSeegert2022Stata}
Bertanha, M., A. H. McCallum, A. Payne, and N. Seegert (2022).
\texttt{bunching}: Stata module to estimate bunching.
\emph{Stata Journal} 22(3), 597--624.

\bibitem[Bosch, Dekker, and Strohmaier(2020)]{BoschDekkerStrohmaier2020Window}
Bosch, N., V. Dekker, and K. Strohmaier (2020).
A data-driven procedure to determine the bunching window: an application to the Netherlands.
\emph{International Tax and Public Finance} 27, 951--979.

\bibitem[Chetty et al.(2011)]{ChettyFriedmanOlsenPistaferri2011Adjustment}
Chetty, R., J. N. Friedman, T. Olsen, and L. Pistaferri (2011).
Adjustment costs, firm responses, and micro vs. macro labor supply elasticities.
\emph{Quarterly Journal of Economics} 126(2), 749--804.

\bibitem[Davezies and Le Barbanchon(2017)]{DaveziesLeBarbanchon2017RDMeasurementError}
Davezies, L. and T. Le Barbanchon (2017).
Regression discontinuity design with continuous measurement error in the running variable.
\emph{Journal of Econometrics} 200(2), 260--281.

\bibitem[Dong(2015)]{Dong2015RoundingRD}
Dong, Y. (2015).
Regression discontinuity applications with rounding errors in the running variable.
\emph{Journal of Applied Econometrics} 30(3), 422--446.

\bibitem[Dong and Kolesar(2023)]{DongKolesar2023IgnoreMeasurementError}
Dong, Y. and M. Kolesar (2023).
When can we ignore measurement error in the running variable?
\emph{Journal of Applied Econometrics} 38(5), 735--750.

\bibitem[Kleven and Waseem(2013)]{KlevenWaseem2013Notches}
Kleven, H. J. and M. Waseem (2013).
Using notches to uncover optimization frictions and structural elasticities.
\emph{Quarterly Journal of Economics} 128(2), 669--723.

\bibitem[Kolesar and Rothe(2018)]{KolesarRothe2018DiscreteRD}
Kolesar, M. and C. Rothe (2018).
Inference in regression discontinuity designs with a discrete running variable.
\emph{American Economic Review} 108(8), 2277--2304.

\bibitem[Manski and Tamer(2002)]{ManskiTamer2002IntervalData}
Manski, C. F. and E. Tamer (2002).
Inference on regressions with interval data on a regressor or outcome.
\emph{Econometrica} 70(2), 519--546.

\bibitem[McCrary(2008)]{McCrary2008DensityTest}
McCrary, J. (2008).
Manipulation of the running variable in the regression discontinuity design: A density test.
\emph{Journal of Econometrics} 142(2), 698--714.

\bibitem[Noack and Rothe(2023)]{NoackRothe2023DonutRD}
Noack, C. and C. Rothe (2023).
Bias and variance of donut regression discontinuity estimators.
arXiv:2308.14464.

\bibitem[Saez(2010)]{Saez2010Bunching}
Saez, E. (2010).
Do taxpayers bunch at kink points?
\emph{American Economic Journal: Economic Policy} 2(3), 180--212.

\bibitem[Song(2024)]{Song2024GeneralBunching}
Song, M. (2024).
Identification and inference for general bunching designs.
arXiv:2411.03625.
\end{thebibliography}

\end{document}
